% Abstract — target 200-250 words for PLOS Comp Bio, up to 350 for Frontiers
% in Comp Neuro. Numbers from results/2026-09-25_ftgrid and round 6.

\begin{abstract}
\noindent
Computational models of the spinal central pattern generator (CPG) in the
Rybak/Danner/Zhang lineage capture interlimb coordination in detail but, by
their authors' own account, omit motoneurons, muscle dynamics, sensory
feedback and synaptic plasticity, and prescribe every synaptic weight by
hand. We asked whether a spinal CPG closed in a loop with its periphery can
tune its own synapses to a working gait. We built a two-leg spiking model of
the rat lumbar CPG in NEST: Izhikevich rhythm-generator half-centres with
asymmetric reciprocal inhibition, motoneuron pools driving muscle proxies with
force, length and fatigue, Ia feedback into the reciprocal-inhibition core, and
spike-timing-dependent plasticity (STDP) on the cutaneous and brainstem inputs
to the rhythm generators. Stance is not ended by a clock: each leg's stance
ends when its own extensor force falls. From ten initial weight distributions
($0$--$16$\,pA) the plastic synapses converge to the same values (cutaneous
$\approx61$\,pA, brainstem $\approx16$\,pA) and produce the same alternating
gait. Across four locomotion modes --- slow, medium and fast walk, and toe
stepping under partial unloading --- and five learning rates, learning that
converges within the session yields force-triggered, alternating stepping in
every mode (extensor--flexor force correlation $-0.47$ to $-0.75$;
left--right $-0.32$ to $-0.81$). Without learning the rhythm falls back on a
safety timer, and a partially learned synapse produces the least coordinated
gait of all. Partial unloading slows learning without abolishing the gait.
The sensorimotor gains of a spinal CPG can thus be learned rather than
prescribed, and whether the closed loop takes over depends on how much
plasticity is available --- a mechanistic account relevant to activity-based
rehabilitation after incomplete spinal cord injury.
\end{abstract}

\noindent\textbf{Keywords:} spinal central pattern generator, STDP,
locomotion, sensory feedback, force-triggered stepping, spinal cord injury,
rehabilitation, NEST simulator.
